\chapter*{Abstract}
\addcontentsline{toc}{chapter}{Abstract}

Answering multi-hop questions over tables requires chaining evidence across table rows and the passages linked from them. Every existing table reasoning pipeline makes this choice of evidence with a fixed rule: embedding similarity, a greedy ranker or a language-model prompt. This project treats the choice as a sequential decision problem. A pointer-network policy walks a per-question graph of rows and hyperlinked passages, selecting one entity at a time or stopping, and a frozen extractive reader answers from the selected evidence. The policy is trained with proximal policy optimisation (PPO) and generalised advantage estimation in an actor--critic setup, on a reward that combines the reader's answer F1, a cost per selected entity and an evidence-recall term that can enter as potential-based shaping or as a terminal bonus.

A review of more than one hundred papers identifies the twelve most relevant and seven research gaps. The central gap is that no reinforcement-learning method selects table entities step by step with a learned stop, and the only RL controller for table--text question answering did no better than a fixed strategy.

On the full HybridQA dev set (3{,}466 questions), with the same frozen reader for every selector, PPO trained from scratch reaches \fScratch{} F1. The best heuristic (BM25 top-3) reaches \fBmThree{}, a HYBRIDER-style greedy rank-then-hop rule \fGreedy{}, and behaviour cloning on the traced evidence \fBc{}; the paired-bootstrap gains of \dScratchBm{} and \dScratchBc{} points have $p<0.001$. The gain holds on tables unseen in training. It also holds for a second reader that was never used in training: on 1{,}000 dev questions PPO from scratch reaches \robScratch{} F1 with RoBERTa-base, against \robBmThree{} for BM25 top-3.
 Warm-starting the policy by behaviour cloning, as the literature recommends, left PPO at the cloned policy's level, and GRPO degraded, most likely because its group normalisation amplified the entity cost. The report describes the formulation, the implementation, the experimental protocol and the plan for the remaining work.
